\section{任意次元のJacobianブロックに対する多変数量子化}
\label{sec:pg-all-dimensions}
本節は三変数一ブロックの普遍Hesse中心公式を、
\emph{任意の次元}の滑らかな三次超曲面に対する
\emph{複数の量子化パラメータを持つ同時量子化}へ拡張する。
論文全体の正規形定理\ref{thm:pg-normal}を入力とするが、
本節での各中心三次元の構成には
形式的中心元の存在補題\ref{lem:pg-formal-center}を使わない。
以下の滑らかな三変数制限の仮定は、
元の超曲面$X$の滑らかさとは異なる追加仮定である。

\subsection{Jacobianブロックの可換Poisson線形族}
$n\ge3$、$N=n+2=3k+m$（$k\ge1,m\ge0$）とし、
\[
 V=\bigoplus_{s=1}^{k}W_s\oplus U,\quad
 \dim W_s=3,\quad\dim U=m,\quad
 F=\sum_{s=1}^k f_s(x_s,u)\in\CC[V]_3.
 \tag{AD}\label{eq:pg-ad-normal}
\]
ここで$F$は滑らかな射影三次超曲面
$X=V(F)\subset\PP^{n+1}$を定め、
\[
 g_s(x_s)=f_s(x_s,0)
 \tag{SM}\label{eq:pg-ad-smooth}
\]
は各$s$について滑らかな平面三次曲線を定めると仮定する。
中心三次部分の$f_s$間での分配は任意に固定する。

$P_s=\partial_{x_{s1}}\wedge\partial_{x_{s2}}
\wedge\partial_{x_{s3}}$、
$q_s=\iota_{df_s}P_s$とおく。
$q_s$は$x_s$方向のみにベクトル添字を持ち、
その係数は$x_s,u$だけに依存する。

\begin{proposition}[一般次元の可換Poisson pencil]
\label{prop:pg-ad-compatible}
各$q_s$は$F$をCasimirとするPoisson双ベクトルであり、
\[
 [q_s,q_t]_{\mathrm{Sch}}=0\qquad(s,t=1,\dots,k).
\]
よって任意の$\boldsymbol\lambda=(\lambda_1,\ldots,\lambda_k)$に対し
\[
 q_{\boldsymbol\lambda}=\sum_s\lambda_s q_s,\qquad
 \pi_{\boldsymbol\lambda}=q_{\boldsymbol\lambda}|_X
\]
はPoisson構造である。
各$q_s|_X$は一次独立であり、
$\pi_{\boldsymbol\lambda}$の一般Poisson階数は
$2\,\#\{s:\lambda_s\ne0\}$となる。
したがって$\boldsymbol\lambda$は$k$次元の線形Poisson族を与える。
\end{proposition}
\begin{proof}
三変数Jacobian括弧のJacobi恒等式は直接成立する。
$F$の$x_s$偏微分は$f_s$のそれに等しいので
$q_s(F,-)=0$。
$s\ne t$なら、$q_s$のベクトル添字は$x_s$方向、
$q_t$の係数は$x_t,u$だけに依存する。
したがって$q_s$の導分は$q_t$の係数を消し、
逆向きも同様なのでSchouten括弧は零である。

$g_s$の滑らかさにより各$q_s$は非零であり、
異なるブロックの$\bigwedge^2W_s$成分に属するから一次独立である。
一般点で各非零係数のブロックは階数二を持ち、
直和なのでアフィン錐上の一般階数は$2r$である
（$r=\#\{s:\lambda_s\ne0\}$）。
射影化で階数が落ちるのはEuler方向が像に入るときだけである。
$m>0$なら$u\ne0$の稠密開集合上でEuler方向は中心成分を持ち、
像はactive方向のみなので階数は落ちない。
$m=0$の場合、$r<k$なら非零のinactive座標成分から同様に落ちず、
$r=k$ならEulerが像に入るためには
$3g_s(x_s)=\sum_i x_{si}g_{s,x_{si}}=0$が全$s$で必要となる。
$X$が$\bigcap_sV(g_s)$に含まれることはない
（$k\ge2$のとき各ブロックは非零三次式である）。
よって一般点では階数は落ちない。
\end{proof}

\subsection{任意階数に対する独立パラメータ量子化}
$\boldsymbol\eta=(\eta_1,\ldots,\eta_k)$を次数零の独立な変数とし、
$R=\CC[u_1,\ldots,u_m]$とおく。
$R[\boldsymbol\eta]$上の代数$\mathcal A_{\boldsymbol\eta}$を
全active生成元$x_{si}$、中心生成元$u_\alpha$、
および関係式
\begin{equation}
\begin{aligned}
 [x_{s2},x_{s3}]
 &=\eta_s\operatorname{sym}_R\bigl((f_s)_{x_{s1}}\bigr),\\
 [x_{s3},x_{s1}]
 &=\eta_s\operatorname{sym}_R\bigl((f_s)_{x_{s2}}\bigr),\\
 [x_{s1},x_{s2}]
 &=\eta_s\operatorname{sym}_R\bigl((f_s)_{x_{s3}}\bigr),\\
 [x_{si},x_{tj}]&=0\ (s\ne t),\qquad [u_\alpha,x_{si}]=0
\end{aligned}
\tag{AQ}\label{eq:pg-ad-relations}
\end{equation}
で定める。すなわち
\[
 \mathcal A_{\boldsymbol\eta}
 =\mathcal A_{f_1,\eta_1}
 \otimes_{R[\boldsymbol\eta]}\cdots
 \otimes_{R[\boldsymbol\eta]}\mathcal A_{f_k,\eta_k}.
 \tag{TP}\label{eq:pg-ad-tensor}
\]
ここでは各ブロックの係数に\emph{同じ中心変数$u$}を使う。
積のテンソルは中心係数環上で取り、
中心変数をブロックごとに複製しない。

\begin{theorem}[任意次元の明示的な同時代数量子化]
\label{thm:pg-all-dimensions}
\eqref{eq:pg-ad-normal}--\eqref{eq:pg-ad-smooth}を仮定する。
各$f_s$に対し、
系\ref{cor:pg-smooth-ternary-explicit}の
有限多項式の中心三次元$\Omega_{s,\eta_s}$を構成する。
すると
\begin{equation}
 \boxed{\qquad
 \Omega_{\boldsymbol\eta}=\sum_{s=1}^k\Omega_{s,\eta_s},
 \qquad
 \Omega_{\boldsymbol0}=F
 \qquad}
 \tag{AC}\label{eq:pg-ad-universal-center}
\end{equation}
は$\mathcal A_{\boldsymbol\eta}$の\emph{厳密な中心元}である。
また原点を含むZariski開集合
$\boldsymbol0\in\mathcal U\subset\mathbb A^k_{\boldsymbol\eta}$を選び、
以下を全て満たせる。
\begin{enumerate}[label=(\arabic*)]
\item 各$\boldsymbol a\in\mathcal U$について
$A_{\boldsymbol a}=\mathcal A_{\boldsymbol\eta}
 \otimes_{\CC[\boldsymbol\eta]}\CC_{\boldsymbol a}$
は$N$生成元、$\binom N2$個の二次関係式による
Koszul AS-regular代数であり、
\[
 \operatorname{gldim}A_{\boldsymbol a}=N,\quad
 \operatorname{Hilb}_{A_{\boldsymbol a}}(z)=(1-z)^{-N},
 \quad\text{AS parameter}=N.
\]
各内部次数成分は$\OO(\mathcal U)$上局所自由である。
\item $\Omega_{\boldsymbol a}$は全てのファイバーで
左右正則元となる。中心三次商
$B_{\boldsymbol a}=A_{\boldsymbol a}/(\Omega_{\boldsymbol a})$
はAS-Gorensteinで
\[
 \operatorname{injdim}B_{\boldsymbol a}=N-1,\quad
 \text{AS parameter}=N-3=n-1,\quad
 \operatorname{Hilb}_{B_{\boldsymbol a}}(z)
 =\frac{1-z^3}{(1-z)^N}.
\]
$\mathcal B=\mathcal A/(\Omega)$の各内部次数成分も
$\OO(\mathcal U)$上局所自由であり、
$\mathcal B_{\boldsymbol0}=\CC[V]/(F)$である。
\item 原点での独立な一次変形方向は
命題\ref{prop:pg-ad-compatible}の
$k$個のPoisson構造$q_s|_X$に対応する。
任意の$\boldsymbol\lambda\in\CC^k$について
$\eta_s=\lambda_s\hbar$と特殊化すれば、
$\pi_{\boldsymbol\lambda}$の
\emph{一変数代数的量子化}が得られる。
その中心三次元も式\eqref{eq:pg-ad-universal-center}で明示される。
\end{enumerate}
Hesse座標を各ブロックごとに選ぶと、
$\Omega_{\boldsymbol\eta}$は各$\eta_s$に関して
次数高々四の\emph{有限多項式}であり、
異なるパラメータ$\eta_s,\eta_t$を同時に含む
新たな補正項は現れない。
したがって$3k+m=n+2$が任意であっても、
この構成は形式的中心元存在補題
\ref{lem:pg-formal-center}を必要としない。
\end{theorem}
\begin{proof}
各$f_s$の純active三次$g_s$が滑らかなので、
Hesse標準形への線形変換と中心変数の平行移動により、
系\ref{cor:pg-smooth-ternary-explicit}を適用できる。
その公式は係数を数値特殊化する必要のない
$\eta_s^2$の高々二次の多項式であり、
$\Omega_{s,0}=f_s$、各active生成元との交換子が零である。
他ブロックの生成元は$x_s$と可換であるから
$\Omega_{s,\eta_s}$は$\mathcal A_{\boldsymbol\eta}$全体で中心となり、
和も中心である。
これで式\eqref{eq:pg-ad-universal-center}を得る。
特に\emph{異なるブロック間の量子的混合補正は必要ない}。

命題\ref{prop:pg-PBW}を各$s$と各$\eta_s$について適用する。
独立なパラメータを持つので、原点を含む各ブロックの
Zariski開集合の積を取ればよい。
$R[\boldsymbol\eta]$-自由なPBW基底を各因子で選び、
\eqref{eq:pg-ad-tensor}でテンソルすると、
全体の各内部次数は底上局所自由で、
$u_1,\ldots,u_m$は各ファイバーで中心正則列である。
その商は
\[
 A_{\boldsymbol a}/(u_1,\ldots,u_m)
 \simeq\bigotimes_{s=1}^k
  S_{s,a_s},\qquad
 S_{s,a_s}=\mathcal A_{f_s,a_s}/(u).
\]
各$S_{s,a_s}$は三次元Koszul AS-regular domain。
それらのKoszul分解と双対のテンソル積から、
右辺は次元$3k$、AS parameter $3k$の
Koszul AS-regular代数となる。
補題\ref{lem:pg-central-linear}を
$m$回適用し、$A_{\boldsymbol a}$の
次元$N=3k+m$のKoszul AS-regular性を得る。
PBW基底がHilbert級数と二次関係数を与える。

$\Omega$の左右正則性には
定理\ref{thm:pg-quantization}の証明と同じ
\emph{一ブロックactive次数}のfiltrationを使う。
第一ブロック$x_1$の次数を数え、
他のactive変数と$u$のfiltration次数を零とする。
$\operatorname{gr}_{x_1}A_{\boldsymbol a}
\simeq S_{1,a_1}\otimes_\CC
 A_{\mathrm{rest},\boldsymbol a}$
（共通中心変数に関して$R$を介したテンソル）
が得られる。
$\Omega_{1,a_1}$のactive次数三の最高記号は、
$\boldsymbol a=\boldsymbol0$で$g_1\ne0$なので、
必要なら開集合をさらに縮めることにより
$S_{1,a_1}$の非零元となる。
他の$\Omega_{s,a_s}$はactive次数零なので、
$\Omega_{\boldsymbol a}$の最高記号は
$\Omega_{1,a_1}$の非零active三次記号である。
$S_{1,a_1}$はdomainであり、
$A_{\mathrm{rest},\boldsymbol a}$は$R$-自由なので、
最高記号の左右乗法は単射。
ゆえに$\Omega_{\boldsymbol a}$は左右正則元である。
補題\ref{lem:pg-central-cubic}で商のAS-Gorenstein性を得る。

各内部次数の$\mathcal A$は底上有限局所自由で、
$\Omega_{\boldsymbol a}$の正則性により
\[
 0\longrightarrow(A_{\boldsymbol a})_{d-3}
 \xrightarrow{\cdot\Omega_{\boldsymbol a}}
 (A_{\boldsymbol a})_d
 \longrightarrow(B_{\boldsymbol a})_d
 \longrightarrow0
\]
のrankは全$\boldsymbol a\in\mathcal U$で一定である。
有限表示行列の定階数条件より、
$\mathcal B_d$も$\OO(\mathcal U)$上局所自由である。
よって多変数底上の平坦性を得る。

最後に定義関係式の$\eta_s$線形項と
$\Omega_{\boldsymbol\eta}\bmod\boldsymbol\eta=F$から、
原点での一次Poisson括弧が$q_s$であることが分かる。
$\eta_s=\lambda_s\hbar$による底変更が
全ての方向を同時に実現する。
\end{proof}

\begin{corollary}[一般次元での全方向の有限Hessian公式]
\label{cor:pg-ad-hessian}
各ブロックのHesse正規化と中心平行移動を済ませ、
その際の量子化変数を$\eta_s$と書く。
$f_s$の混合係数を$A_s,B_s,C_s,D_s,E_s,H_s,G_s$、
Hesse係数を$t_s$とする。
\[
 {\mathscr H}_s=
 \det\left(\frac{\partial^2 f_s}
 {\partial x_{si}\partial x_{sj}}\right)_{i,j=1}^3.
\]
系\ref{cor:pg-hesse-hessian-form}の
有限多項式$\Psi_{t_s}(f_s)$を用いると、
\[
 \boxed{\quad
 \Omega_{\boldsymbol\eta}=
 \operatorname{sym}\left(
 F-\sum_{s=1}^k\frac{\eta_s^2}{12}{\mathscr H}_s
 +\sum_{s=1}^k\frac{\eta_s^2}{6}A_sB_sC_s
 +\sum_{s=1}^k\eta_s^4\Psi_{t_s}(f_s)
 \right).
 \quad}
 \tag{AH}\label{eq:pg-ad-hessian-center}
\]
これは$\boldsymbol\eta$の多項式であり、
全次数で厳密な中心元である。
特に$\eta_s=\lambda_s\hbar$としたとき
$\hbar^2$のHessian補正は$\lambda_s^2$、
$\hbar^4$補正は$\lambda_s^4$を係数とする。
異なるブロックの混合Hessian項
$\eta_s^2\eta_t^2$は生じない。
一般の$g_s$へはHesse正規化時の
$\eta_s$の行列式スケール変換を施してから適用する。
\end{corollary}
\begin{proof}
系\ref{cor:pg-hesse-hessian-form}の式を各ブロックに用いて
和を取れば得られる。
各補正は$x_s,u$にのみ依存し、
他ブロックの生成元と可換なため、
高次の交叉補正を必要としない。
\end{proof}

\begin{remark}[全次元・全成分についての正確な範囲]
\label{rem:pg-ad-scope}
$k$は$1\le k\le\lfloor(n+2)/3\rfloor$を動ける。
この上限は正規形における\emph{許される最大ブロック数}であり、
全Poissonモジュライの成分数や族の次元と同一視しない。
固定したブロック分解における三次式$F$の線形パラメータ空間では、
滑らかな$X$と全ての滑らかな$g_s$という条件は
非空なZariski開条件である。
実際$\sum_{s,i}x_{si}^3+\sum_\alpha u_\alpha^3$
はその一例である。
従って本定理は任意次元の各正規形型の
\emph{非空な一般部分}に適用される。
ただし、滑らかな$X$であっても一部の$g_s$が
特異な場合はあり、その部分への無条件な適用は主張しない。

固定した$F$と固定したブロック分解に対する
$k$変数量子化を構成したのであり、
相対Poissonモジュライの\emph{全既約成分}の
普遍的量子化空間を構成したという意味ではない。
またこの量子化から$\operatorname{qgr}B$の
幾何学的smoothnessや可換Kuznetsov復元は従わない。
Hesse標準形を全係数空間にわたって
正準的に同時選択できるとも主張しない。
\end{remark}
